\section{总结与延伸}

如果把整门 CS231N 看成一条主线，那么 Lecture 18 做的是一次``高层反编译''：它把前面学过的卷积网络、识别、检测、分割、生成、多模态与具身智能问题，重新还原成三个更大的研究命题。

\subsection{命题一：视觉 AI 首先是对世界结构的学习}

从进化视角到对象识别史，这一讲提醒我们，视觉不是随机输入，而是世界结构最密集的表达之一。对象、属性、关系、语言、动作，这些层层递进的任务共同构成了视觉智能的扩展路径。

\subsection{命题二：AI 的价值不只在于复制人类，还在于补足人类}

无论是 fine-grained recognition、社会统计、医疗监测，还是长期持续感知任务，AI 都可以成为一种超越人类尺度、稳定性和覆盖面的工具。但这类能力只有在公平、隐私与责任边界被认真对待时，才是真正有价值的进步。

\subsection{命题三：最难的问题不是能不能做，而是该不该做、该怎么做}

这一讲最终把问题交还给 human-centered AI：自动化什么、保留什么、增强谁、如何评估``有用''，这些都不是模型精度曲线自己会回答的。它们需要研究者主动把人的价值、偏好与制度约束嵌入系统设计。

\subsection{把三条主线翻译成真实系统设计问题}

如果把 lecture18 的三条主线落到工程上，可以得到下面这张很实用的检查表。它的价值在于：\textbf{每做一个视觉系统，都可以拿这张表过一遍，判断自己究竟在解决哪一层问题。}

\begin{table}[H]
\centering
\small
\renewcommand{\arraystretch}{1.18}
\resizebox{\textwidth}{!}{%
\begin{tabular}{p{2.8cm}p{3.6cm}p{3.3cm}p{3.6cm}}
\toprule
\textbf{主线} & \textbf{系统目标} & \textbf{常见失败模式} & \textbf{研究者必须补问的问题} \\
\midrule
See what humans see & 学会可靠的对象、关系、语言与时序理解 & 只会背训练分布，离开 benchmark 就不稳 & 我们究竟学到了世界结构，还是学到了数据集套路？ \\
See what humans don't see & 补足人类在尺度、一致性、疲劳和长尾上的短板 & 把偏差、不公平或监控一起规模化放大 & 模型补的是人的盲点，还是放大了人的弱点？ \\
See what humans want to see & 让能力真正服务人的目标、偏好和责任边界 & 技术上可行，但任务本身不值得自动化 & 谁定义任务？谁承担后果？谁有权拒绝自动化？ \\
\bottomrule
\end{tabular}
}
\caption{lecture18 的三层问题可以直接转成系统设计审查表}
\end{table}

\begin{knowledgebox}{为什么这张表值得记下来}
很多坏系统并不是因为模型不够强，而是因为它们从一开始就答错了问题。它们也许能 ``see''，但没有想清楚 ``why''、``for whom'' 与 ``at what cost''。
\end{knowledgebox}

\subsection{给做视觉研究和做系统的人各留一个提醒}

对研究者来说，这一讲最重要的提醒是：\textbf{不要把 benchmark 上的成功直接误认为世界中的成功。} ImageNet 时代证明了标准化评测的力量，但 lecture18 又反过来提醒，真实世界中的价值判断、隐私边界、长程行为与具身执行，并不会因为 top-1 accuracy 提升而自动解决。

对系统工程师来说，这一讲最重要的提醒是：\textbf{不要把 human-in-the-loop 当成免责条款。} 如果系统设计本身忽视了人的疲劳、注意力上限、偏好差异与责任链条，那么把最终确认塞给人类，只是在把系统缺陷转嫁给一线使用者。

\begin{importantbox}{三条最终 takeaway}
\begin{enumerate}
    \item \textbf{视觉 AI 的历史，是把人类感知能力转化为可学习系统的历史。}
    \item \textbf{现代 AI 最有价值的方向之一，是补足人类盲点，而不是机械复制人类。}
    \item \textbf{能力越强，越需要 human perspective；否则 AI 放大的不只是效率，也可能是不平等、监控与错误。}
\end{enumerate}
\end{importantbox}

把这一讲放在整门课最后，其实非常合适。因为学完模型之后，学生最容易产生的错觉是：只要模型更大、数据更多、benchmark 更高，问题就自然解决了。Lecture 18 明确告诉我们并非如此。真正成熟的视觉研究者，需要同时理解\textbf{表示、应用、风险与价值}。

\begin{warningbox}{这也是整门课最值得带走的一句话}
我们当然要 build AI to see things，但更重要的是 build AI to help people, respect people, and remain answerable to people.
\end{warningbox}

\subsection{本章小结}
Lecture 18 的真正价值，不是再教一个新模型，而是把整门 CS231N 的技术积累重新放回 ``人如何使用视觉智能'' 这个更大的问题里。它要求我们同时理解世界结构、能力边界、风险治理与人的真实偏好，这也是现代视觉系统从论文走向现实部署时最难跳过的一关。

\subsection{一页总表：从课程主线到系统设计}

\begin{table}[H]
\centering
\small
\renewcommand{\arraystretch}{1.18}
\begin{tabular}{p{2.8cm}p{3.4cm}p{3.6cm}p{4.0cm}}
\toprule
\textbf{视角} & \textbf{课程里回答的问题} & \textbf{落到系统中的表现} & \textbf{设计时最该追问什么} \\
\midrule
感知建模 & 模型是否真正学到对象、关系与时序结构 & 检测、分割、VLM、具身感知等能力是否稳健 & 系统是在理解环境，还是在复述数据集偏差？ \\
能力增强 & AI 是否补足人类看不到或难以稳定完成的部分 & 医疗监测、社会统计、长期感知、辅助技术 & 我们是在扩展人的能力，还是放大错误与不公平？ \\
人本约束 & 自动化边界由谁定义、如何被审查 & 任务选择、反馈环路、责任划分与拒绝机制 & 谁定义目标、谁承担后果、谁能说不？ \\
\bottomrule
\end{tabular}
\caption{把 Lecture 18 的三层主线压缩成系统设计检查表}
\end{table}

\subsection{给继续学习视觉 AI 的读者三条延伸路线}

\begin{importantbox}{如果把 Lecture 18 当成分叉点，后续最值得追的三条线}
\begin{enumerate}
    \item \textbf{具身与机器人视觉}：继续往 action、planning、real-to-sim transfer 与数字孪生走，理解视觉如何进入闭环控制系统；
    \item \textbf{医疗与辅助技术}：关注 AI 在诊断、长期监测、残障辅助与脑机接口中的真实价值与责任边界；
    \item \textbf{Human-centered evaluation}：不只看 benchmark，而是研究人类偏好、制度约束与部署后的社会反馈。
\end{enumerate}
\end{importantbox}

\begin{knowledgebox}{为什么这三条路线比 ``再刷一个 benchmark'' 更重要}
Lecture 18 的立场非常明确：视觉研究真正走向成熟，不是因为模型更大了，而是因为它开始正面回答 ``谁受益''、``谁承担代价''、``什么任务值得自动化'' 这些问题。未来越是强大的视觉系统，越需要同时接受技术评估和人本评估。
\end{knowledgebox}

\subsection{拓展阅读}

\begin{itemize}
    \item Human-centered AI 与 value-sensitive design 相关文献，帮助把 ``用户偏好'' 变成可执行的系统设计原则
    \item 具身智能与机器人学习课程，补上从视觉感知到行动闭环的部分
    \item AI fairness、privacy、accountability 方向的论文与案例研究，理解视觉系统部署后的长期后果
    \item 医疗 AI 与 assistive robotics 论文，观察 ``see what humans don't see'' 如何在真实场景中创造价值
\end{itemize}

\end{document}
